\chapter{Implementing Cryptography Safely}

Every chapter so far taught you how a cryptographic machine works on the inside.
You can now trace AES through its rounds, run RSA with small primes by hand, see
why GCM dies on a repeated nonce, and follow a TLS 1.3 handshake message by
message. That knowledge is real and it matters. But there is a gap between
understanding a primitive and shipping a system that uses it, and that gap is
where almost all real-world cryptographic failures live. The math is rarely what
breaks. What breaks is the wiring: a nonce that repeats because a counter reset, a
key that lived in a log file, a comparison that returned early and leaked a tag
byte at a time, a downgrade that nobody tested, a ``temporary'' ECB mode that
shipped to production.

The conventional advice ``do not roll your own crypto'' highlights a genuine risk,
but provides little constructive guidance. The operational rule is more specific:
\emph{do not implement primitives, do compose audited high-level
constructions, and treat every place where you must make a choice as a place where
you can make a mistake.} The craft of safe implementation is the discipline of
removing your own choices --- nonces, modes, padding, key sizes, comparison
routines --- and handing them to code that has already been attacked by people
smarter and meaner than your future self at 3 a.m.

We will work through the library hierarchy (why libsodium and Tink exist and what
``misuse-resistant'' actually buys you), key management and rotation with a worked
key hierarchy, secure defaults, testing against published test vectors, a real
threat-modeling pass on a small system, the genuinely rare cases where you do
touch a primitive, and a one-page checklist you can paste into a code review.
There is less new mathematics here than anywhere else in the book and more
hard-won engineering judgment. Both keep secrets. Only one of them is usually the
reason your secret leaked.

\chapterline

\section{Defining the Core Rule}
\label{ch39:one-rule}

Let us start with the slogan and dismantle it carefully, because half of using it
well is knowing exactly what it forbids.

\begin{focusbox}
\textbf{``Don't roll your own crypto'' decomposed.} The slogan bundles three very
different activities that deserve different verdicts.
\begin{enumerate}
  \item \textbf{Designing a new primitive} --- a new cipher, hash, or signature
  scheme. For almost everyone, almost always: \emph{never}. New primitives need
  years of public cryptanalysis before anyone should trust them. This is a
  research activity, not an engineering one.
  \item \textbf{Implementing a standard primitive} --- your own AES, your own
  Curve25519. \emph{Almost never.} The algorithm is public and correct on paper,
  but a correct-looking implementation can still leak the key through timing,
  cache behavior, or a faulty bignum routine. We will see when --- rarely --- this
  is justified.
  \item \textbf{Composing primitives into a protocol} --- gluing key exchange,
  encryption, and authentication into a secure channel. This is the dangerous one
  precisely because it looks like ordinary programming. Most ``roll your own''
  disasters live here: people who would never write their own AES happily invent
  their own handshake, their own MAC-then-encrypt order, their own session
  resumption. \emph{Use a vetted construction instead.}
\end{enumerate}
The corrected rule: \textbf{compose audited high-level constructions; do not
design or implement primitives; and minimize the choices you are allowed to make
at every layer.}
\end{focusbox}

Why is this so much harder than normal software? In ordinary engineering, a bug
announces itself. The button does not click, the total comes out wrong, the test
goes red. Cryptographic bugs are silent. A system that uses a fixed nonce, or that
never actually verifies a signature, or that compares MACs with a short-circuiting
string equality, \emph{works perfectly} in every functional test. Encrypt,
decrypt, round-trip succeeds, ship it. The defect is invisible until an adversary
who is looking for exactly that defect finds it, and by then the failure is a
breach, not a failed unit test. The feedback loop that keeps normal code roughly
correct is missing here. That single fact --- correctness is not testable by the
people writing the code --- is the primary reason to lean on artifacts that
hostile experts have already tested.

\begin{intuitionbox}
Think of a bridge. Nobody pours their own steel or invents their own bolt; they
use rated parts and assemble them by a building code that encodes every past
collapse. Crypto standards are that code: vetted parts (AES, X25519, HMAC) plus
vetted ways of joining them (TLS, AEAD, HKDF), each written after someone found
out the hard way what happens otherwise.
\end{intuitionbox}

\begin{historybox}
The line ``anyone can invent a cipher they cannot break themselves'' is folk
wisdom usually traced to Bruce Schneier's \emph{Memo to the Amateur Cipher
Designer} (1998), and the underlying point is older --- it is Auguste Kerckhoffs's
principle from 1883 turned into career advice. The history of broken homebrew
crypto is long and consistent: 1999's 802.11 WEP paired RC4 with a
non-cryptographic CRC-32 checksum and 24-bit nonces that inevitably collided;
Adobe's 2013 password ``encryption'' used ECB mode with unsalted password hints
sitting right beside it; the MS-CHAPv2 handshake and countless custom
``obfuscation'' layers in mobile apps fall in an afternoon. In response, Daniel J.
Bernstein, Tanja Lange, and Peter Schwabe released NaCl in 2008 to pioneer
``misuse-resistant'' APIs that remove configuration knobs, and Google followed
with Tink in 2017. The pattern in field failures is almost never ``they used a
weak cipher.'' It is almost always ``they used a fine cipher and wired it up wrong.''
\end{historybox}

\section{The Library Hierarchy: Where to Stand}
\label{ch39:hierarchy}

Cryptographic APIs come in layers, and the single most consequential decision you
make is \emph{which layer you build on}. Higher is safer, because higher means
fewer decisions left to you. Picture a ladder (Figure~\ref{ch39:fig-hierarchy}); you want to stand as near the top
as your problem allows.

\begin{figure}[H]
\centering
\begin{tikzpicture}[font=\footnotesize,
  layer/.style={draw, rounded corners=2pt, minimum width=9.6cm, minimum height=0.95cm, align=center, inner sep=3pt}
]
\node[layer, draw=targetgreen!80, fill=targetgreen!15] (l4) at (0, 4.4) {
  \textbf{Layer 4 --- High-Level Protocols (Zero Assembly)} \\
  \scriptsize TLS 1.3, WireGuard, Signal protocol, Cloud KMS. ``Send secret to recipient.''
};
\node[layer, draw=sourceblue!80, fill=sourceblue!12] (l3) at (0, 3.1) {
  \textbf{Layer 3 --- Misuse-Resistant Box APIs (Knobs Removed)} \\
  \scriptsize libsodium \texttt{crypto\_secretbox}, Google Tink \texttt{Aead}. Nonces automatic.
};
\node[layer, draw=quotientorange!80, fill=quotientorange!15] (l2) at (0, 1.8) {
  \textbf{Layer 2 --- Named AEAD Constructions (Manual State)} \\
  \scriptsize AES-GCM, ChaCha20-Poly1305. Caller must guarantee unique nonce.
};
\node[layer, draw=evered!80, fill=evered!12] (l1) at (0, 0.5) {
  \textbf{Layer 1 --- Raw Primitives \& Modes (Manual Composition)} \\
  \scriptsize Raw AES, CBC mode, HMAC. Caller owns IV, padding, MAC order, timing.
};
\node[layer, draw=trapred!90, fill=trapred!20] (l0) at (0, -0.8) {
  \textbf{Layer 0 --- Custom Mathematical Primitives (Roll Your Own)} \\
  \scriptsize Custom S-boxes, bignum routines, custom curves. Silent flaws guaranteed.
};

\draw[-{Stealth[length=3mm]}, very thick, structurepurple] (5.15, -0.8) -- node[right=2pt, font=\scriptsize\bfseries, structurepurple, align=center]{Safer\\(Fewer\\Decisions)} (5.15, 4.4);
\draw[-{Stealth[length=3mm]}, very thick, trapred] (-5.15, 4.4) -- node[left=2pt, font=\scriptsize\bfseries, trapred, align=center]{Hazardous\\(Expanding\\Surface)} (-5.15, -0.8);
\end{tikzpicture}
\caption{The cryptographic API ladder. As you ascend from Layer 0 to Layer 4, the library assumes responsibility for error-prone mechanics (constant-time branches, IV generation, padding, and MAC ordering), narrowing your attack surface to near zero. Every parameter eliminated from an API signature is an exploit vector eliminated from production.}
\label{ch39:fig-hierarchy}
\end{figure}

The mistakes scale with the layer. At Layer 4 there is almost nothing left to get
wrong; you call code whose job is to be secure end to end. At Layer 3 you choose a
key and pass plaintext; the API picks the nonce, the mode, the tag, and refuses to
let you reuse a nonce. At Layer 2 you have re-acquired one sharp edge --- the nonce
--- and GCM punishes a repeat catastrophically (recall Chapter~16). At Layer 1 every sharp edge is back: you choose the mode, generate the
IV, add the padding, decide MAC-then-encrypt versus encrypt-then-MAC, and write the
tag comparison. Layer 0 adds the algorithm itself. \textbf{Each step down the
ladder multiplies the number of independent ways to be silently wrong.}

\begin{intuitionbox}
Read the ladder this way: every API parameter you are \emph{required} to fill in is
a place the library has delegated a security decision to you. Count the parameters.
\texttt{secretbox(message, key)} delegates almost nothing.
\texttt{aes\_gcm\_encrypt(message, key, nonce, aad)} delegates the nonce ---
exactly the parameter with the famous catastrophic failure mode. ``How many ways
can I hold this wrong?'' is, to a first approximation, ``how many arguments does it
take?''
\end{intuitionbox}

\subsection{NaCl and libsodium: the box that hides the nonce}

NaCl (``salt,'' the Networking and Cryptography Library) was Daniel J. Bernstein,
Tanja Lange, and Peter Schwabe's deliberate answer to the misuse problem.
libsodium is its portable, packaged, widely audited descendant. Its philosophy is
this whole chapter compressed into an API: \emph{pick good primitives once, expose
them through operations that are hard to misuse, and delete the knobs.} There is no
cipher selection, no mode menu, no negotiation. Secret-key authenticated encryption
is a single function.

\begin{implementbox}
\textbf{Learning only --- this is illustrative, but here the lesson \emph{is} to
use the real library.} The point of the snippet is how little you are asked to
decide, not to teach you to reimplement anything.
\begin{lstlisting}[language=Python, caption={libsodium-style authenticated
encryption via PyNaCl. Notice how few decisions you are handed.}]
import nacl.secret, nacl.utils

# A 32-byte key from a CSPRNG. Never type a key in; never derive it from
# a password without a real KDF (see the key-management section).
key = nacl.utils.random(nacl.secret.SecretBox.KEY_SIZE)
box = nacl.secret.SecretBox(key)

# Encrypt. We do NOT pass a nonce: SecretBox draws a fresh random one and
# prepends it to the ciphertext. One catastrophic decision removed from us.
ct = box.encrypt(b"transfer 100 to account 7")

# Decrypt. The nonce rides along with the ciphertext; the tag is checked.
# If a single bit was flipped, this RAISES -- it never returns garbage.
pt = box.decrypt(ct)
assert pt == b"transfer 100 to account 7"
\end{lstlisting}
\end{implementbox}

Look at what is \emph{absent}. You did not pick a cipher (it is
XSalsa20-Poly1305). You did not pick a mode. You did not choose between encrypt-then-MAC
and the alternatives --- the construction is fixed and correct. You did not
generate a nonce, so you cannot reuse one. The decryption either returns the exact
plaintext or throws; there is no code path that hands you unauthenticated bytes to
``check later.'' Every one of those absences corresponds to a class of CVEs that
simply cannot occur in code written against this API. That is what
\emph{misuse-resistant} means: not ``hard to break if used correctly,'' but ``hard
to use incorrectly in the first place.''

\begin{focusbox}
\textbf{Why XSalsa20 and not raw AES-GCM here?} The leading \texttt{X} means
\emph{extended nonce} --- 192 bits instead of 96. With a 192-bit random nonce you
can pick nonces at random forever and the chance of a collision stays negligible
(a birthday bound over $2^{192}$; see the connection box). That is precisely what
lets \texttt{SecretBox} generate nonces randomly without bookkeeping. AES-GCM's
96-bit nonce is too short for that, which is why AEAD-at-Layer-2 forces you to
manage a counter. The library chose a primitive whose parameters make the safe
usage the only usage. Design choices propagate.
\end{focusbox}

\begin{connectionbox}
The random-nonce argument is a birthday-bound calculation from the companion
probability book. With an $n$-bit nonce drawn uniformly, the expected number of
draws before a collision is about $\sqrt{2^{n}} = 2^{n/2}$. For $n = 96$ that is
$2^{48}$ messages --- reachable by a busy service, which is why GCM nonces are
managed, not random. For $n = 192$ it is $2^{96}$ --- astronomically out of reach,
which is why XSalsa20's extended nonce can be random. The whole API ergonomics
question reduces to one square-root.
\end{connectionbox}

\subsection{Tink: types that refuse to be confused}

Google's Tink (in Java, C++, Go, Python) attacks the same problem from the
type-system side. Its central ideas are worth naming because they generalize
beyond any one library.

\textbf{Keys are objects, not byte strings.} In Tink a key is a \texttt{Keyset}
that carries its own algorithm, parameters, and key-id metadata. You cannot
accidentally feed an HMAC key into AES-GCM, because the type that comes out of the
keyset \emph{is} an \texttt{Aead} or a \texttt{Mac}, never raw bytes. A whole genus
of bugs --- ``key confusion,'' where the same key material gets used for two
purposes --- is unrepresentable.

\textbf{Operations are interfaces, not algorithms.} You ask for an \texttt{Aead}
and call \texttt{encrypt(plaintext, associatedData)}. You did not name AES-GCM; the
keyset did. Swapping AES-GCM for ChaCha20-Poly1305, or rotating to a new key, is a
keyset change with no call-site edits. Crypto-agility without per-call decisions.

\textbf{Keysets carry rotation built in.} A keyset can hold several keys at once,
one marked \emph{primary}. New data is encrypted under the primary; incoming data
is decrypted under whichever key-id its header names. Rotation becomes ``add a key,
promote it to primary, later retire the old one'' --- and old ciphertext keeps
decrypting the entire time. We will see why that staging matters when we work
through rotation by hand.

\begin{intuitionbox}
libsodium removes decisions by deleting knobs; Tink removes decisions by making the
wrong combination a type error your compiler rejects before the code runs. Same
goal, two routes: \emph{the safe path is the path of least resistance.} If your
crypto API makes you stop and think hard about nonces and modes on every call, it
is, by this standard, a badly designed API --- and you are one tired afternoon away
from a vulnerability.
\end{intuitionbox}

\section{Misuse-Resistance Made Concrete}
\label{ch39:misuse}

It is worth seeing, with small numbers, exactly which mistake each design choice
forecloses. We will take the single most expensive misuse in modern symmetric
crypto --- nonce reuse in a counter-based AEAD --- and watch a misuse-resistant API
prevent it.

Recall from Chapter~16 that GCM (and any CTR-based
scheme) generates a keystream as a function of the key and nonce alone. Encrypt two
messages $m_1, m_2$ under the same key and \emph{same nonce} and you get keystream
$S$ used twice:
\[
c_1 = m_1 \xor S, \qquad c_2 = m_2 \xor S
\quad\Longrightarrow\quad
c_1 \xor c_2 = m_1 \xor m_2 .
\]
The key vanished; an eavesdropper learns the XOR of two plaintexts, and for GCM
specifically the authentication key is also exposed, letting the attacker forge.
Let us make the confidentiality half tangible with one byte.

\begin{example}[Nonce reuse leaks plaintext, by hand]
\label{ch39:nonce-reuse-ex}
Work in single bytes. Suppose the keystream byte for a given (key, nonce) is
$S = \texttt{0xA3}$. Message one is the ASCII letter \texttt{'Y'} $= \texttt{0x59}$;
message two is \texttt{'N'} $= \texttt{0x4E}$ (a yes/no vote, say). Reusing the
nonce:
\[
c_1 = \texttt{0x59} \xor \texttt{0xA3} = \texttt{0xFA}, \qquad
c_2 = \texttt{0x4E} \xor \texttt{0xA3} = \texttt{0xED}.
\]
Eve never learns $S$. She computes
\[
c_1 \xor c_2 = \texttt{0xFA} \xor \texttt{0xED} = \texttt{0x17}
            = \texttt{0x59} \xor \texttt{0x4E} = \text{`Y'} \xor \text{`N'} .
\]
\texttt{0x17} is the XOR of the two votes. Given that the plaintexts are one of two
known ASCII letters, that single value tells Eve the two votes differed --- and a
sliver of known plaintext (a fixed header, a guessed word) peels the rest apart.
One repeated nonce, and the ``encrypted'' channel is a crib. Now imagine the
messages are kilobytes of structured English: classical two-time-pad cryptanalysis
recovers both in full.
\end{example}

Here is the point of the whole section. At Layer 2 (raw AES-GCM) this disaster is
one off-by-one counter away --- a process restart that resets a nonce counter to
zero, a forked worker that inherits the same counter, a virtual-machine snapshot
restored twice. At Layer 3, the API never lets you supply the nonce, so the bug is
\emph{unreachable from your code}. You did not get more careful; the design got
the carefulness for you.

\begin{trapbox}
The nastiest version of this bug is the VM-snapshot or fork: your nonce counter is
perfectly correct, monotone, never reused --- within one process. Then the deploy
system clones the running machine, or your worker pool forks after seeding the
counter, and two ``different'' processes march the same counter from the same
start. Every functional test passes. The collision happens only in production,
only under scale, only once it matters. This is exactly the class of failure that
a 192-bit random nonce (XSalsa20) or a managed-keyset design sidesteps by
construction.
\end{trapbox}

\section{Key Management: The Part Nobody Tests}
\label{ch39:keymgmt}

A cipher with a key is a machine with a secret. If the secret is mishandled, the
machine's mathematical strength is irrelevant. Key management is where most
real-world crypto actually lives and dies, and it is the part with the fewest unit
tests, because ``the key never leaked'' is not something a test suite can assert.

\subsection{The key lifecycle}

Every key passes through stages, and each stage has its own failure mode.

\begin{center}
\begin{tikzpicture}[font=\footnotesize, node distance=5mm and 5mm,
  st/.style={draw=sourceblue!80, rounded corners=2pt, fill=sourceblue!8, minimum height=0.85cm, align=center, inner sep=3.5pt},
  ar/.style={-{Stealth[length=2.2mm]}, thick, sourceblue!80!black}]
\node[st] (gen) {\textbf{Generate} \\ \scriptsize CSPRNG, full entropy};
\node[st, right=of gen] (store) {\textbf{Store} \\ \scriptsize KMS / HSM / sealed};
\node[st, right=of store] (use) {\textbf{Use} \\ \scriptsize in memory, minimal};
\node[st, right=of use] (rot) {\textbf{Rotate} \\ \scriptsize new primary};
\node[st, right=of rot] (rev) {\textbf{Retire / Revoke} \\ \scriptsize wipe, never reuse};
\draw[ar] (gen)--(store); \draw[ar] (store)--(use); \draw[ar] (use)--(rot);
\draw[ar] (rot)--(rev);
\draw[ar] (rot.south) to[out=-125,in=-55] (use.south);
\end{tikzpicture}
\end{center}

\textbf{Generate} from a cryptographically secure RNG with full entropy ---
\texttt{/dev/urandom}, \texttt{getrandom()}, the OS CSPRNG. Never from
\texttt{rand()}, never from a timestamp, never a hand-typed string. \textbf{Store}
so the key is not sitting in a config file, an environment variable dumped to logs,
or a git repository. \textbf{Use} for the shortest time and smallest scope
possible, ideally never copying the raw key out of a hardware security module.
\textbf{Rotate} on a schedule and after any suspected exposure. \textbf{Retire} by
wiping the material and never, ever reusing it for a new purpose.

\begin{trapbox}
The most common key-management failure is the most boring one: the key in source
control. A long-lived API key or signing key committed to git, then the repo goes
public or a laptop is stolen, and the key is in history forever --- removing the
file does not remove it from the commit graph. The fix is procedural, not
mathematical: secrets never touch the repo, they come from a secrets manager or
environment injected at deploy time, and a pre-commit scanner (gitleaks,
trufflehog) blocks the commit. No cipher protects you from publishing the key.
\end{trapbox}

\subsection{Key hierarchies and key wrapping}

You almost never encrypt data directly under a long-lived master key. Instead you
build a \emph{hierarchy}: one well-protected root, many cheap leaves. This is
\emph{envelope encryption}, and it solves several problems at once.

The pattern: a \textbf{master key} (also called a key-encryption key, KEK) lives in
an HSM or a managed KMS and never leaves it. For each chunk of data --- a file, a
record, a session --- you generate a fresh \textbf{data-encryption key} (DEK),
encrypt the data with the DEK, then ask the KMS to encrypt (``wrap'') the DEK under
the KEK. You store the ciphertext and the wrapped DEK together; the bare DEK is
wiped from memory. To read, you send the wrapped DEK to the KMS, it unwraps with
the KEK, and you decrypt (Figure~\ref{ch39:fig-envelope}).

\begin{figure}[H]
\centering
\begin{tikzpicture}
  % Zone 1: KMS / Hardware Security Module Boundary
  \draw[draw=structurepurple!60, fill=structurepurple!5, rounded corners=4pt, dashed, thick] (-0.3, 4.1) rectangle (12.3, 6.4);
  \node[font=\scriptsize\bfseries, structurepurple, anchor=west] at (0.0, 6.1) {ZONE 1: KMS / HSM Trust Boundary (Hardware Root of Trust)};

  \node[draw=structurepurple!80, fill=structurepurple!15, rounded corners=2pt, font=\scriptsize\bfseries, align=center, minimum width=4.5cm, minimum height=0.9cm] (kek) at (6.0, 5.0) {
    Master Key (KEK) \\ \tiny Held in tamper-resistant HSM. Never exported to application memory.
  };

  % Zone 2: Application Process Memory Boundary
  \draw[draw=sourceblue!60, fill=sourceblue!5, rounded corners=4pt, dashed, thick] (-0.3, 0.9) rectangle (12.3, 3.2);
  \node[font=\scriptsize\bfseries, sourceblue!80!black, anchor=west] at (0.0, 2.9) {ZONE 2: Application Process Runtime Memory};

  \node[draw=sourceblue!80, fill=sourceblue!15, rounded corners=2pt, font=\scriptsize\bfseries, align=center, minimum width=3.2cm] (dek1) at (2.0, 1.9) {
    DEK$_1$ (Ephemeral) \\ \tiny Fresh 256-bit key from CSPRNG
  };
  \node[draw=sourceblue!80, fill=sourceblue!15, rounded corners=2pt, font=\scriptsize\bfseries, align=center, minimum width=3.2cm] (dek2) at (6.0, 1.9) {
    DEK$_2$ (Ephemeral) \\ \tiny Wiped via \texttt{explicit\_bzero}
  };
  \node[draw=sourceblue!80, fill=sourceblue!15, rounded corners=2pt, font=\scriptsize\bfseries, align=center, minimum width=3.2cm] (dek3) at (10.0, 1.9) {
    DEK$_3$ (Ephemeral) \\ \tiny One key per file/record
  };

  % Zone 3: Persistent Storage (Untrusted)
  \draw[draw=targetgreen!60, fill=targetgreen!5, rounded corners=4pt, dashed, thick] (-0.3, -2.3) rectangle (12.3, 0.1);
  \node[font=\scriptsize\bfseries, targetgreen!80!black, anchor=west] at (0.0, -0.2) {ZONE 3: Persistent Untrusted Storage (Database / S3 / Disk)};

  \node[draw=targetgreen!80, fill=targetgreen!15, rounded corners=2pt, font=\scriptsize, align=center, minimum width=3.4cm] (f1) at (2.0, -1.25) {
    \textbf{Record 1 Payload} \\ \tiny Ciphertext $C_1 = \Enc_{\text{DEK}_1}(M_1)$ \\ \tiny Wrapped Key: $\Enc_{\text{KEK}}(\text{DEK}_1)$
  };
  \node[draw=targetgreen!80, fill=targetgreen!15, rounded corners=2pt, font=\scriptsize, align=center, minimum width=3.4cm] (f2) at (6.0, -1.25) {
    \textbf{Record 2 Payload} \\ \tiny Ciphertext $C_2 = \Enc_{\text{DEK}_2}(M_2)$ \\ \tiny Wrapped Key: $\Enc_{\text{KEK}}(\text{DEK}_2)$
  };
  \node[draw=targetgreen!80, fill=targetgreen!15, rounded corners=2pt, font=\scriptsize, align=center, minimum width=3.4cm] (f3) at (10.0, -1.25) {
    \textbf{Record 3 Payload} \\ \tiny Ciphertext $C_3 = \Enc_{\text{DEK}_3}(M_3)$ \\ \tiny Wrapped Key: $\Enc_{\text{KEK}}(\text{DEK}_3)$
  };

  % Flow arrows
  \draw[-{Stealth[length=2.5mm]}, thick, structurepurple] (kek.south) -- node[pos=0.42, above left=0pt and 2pt, font=\tiny\bfseries, fill=white, inner sep=1pt] {Wraps} (dek1.north);
  \draw[-{Stealth[length=2.5mm]}, thick, structurepurple] (kek.south) -- node[pos=0.42, right=2pt, font=\tiny\bfseries, fill=white, inner sep=1pt] {Wraps} (dek2.north);
  \draw[-{Stealth[length=2.5mm]}, thick, structurepurple] (kek.south) -- node[pos=0.42, above right=0pt and 2pt, font=\tiny\bfseries, fill=white, inner sep=1pt] {Wraps} (dek3.north);

  \draw[-{Stealth[length=2.5mm]}, thick, sourceblue] (dek1.south) -- node[pos=0.45, right=2pt, font=\tiny, fill=white, inner sep=1pt] {Encrypts} (f1.north);
  \draw[-{Stealth[length=2.5mm]}, thick, sourceblue] (dek2.south) -- node[pos=0.45, right=2pt, font=\tiny, fill=white, inner sep=1pt] {Encrypts} (f2.north);
  \draw[-{Stealth[length=2.5mm]}, thick, sourceblue] (dek3.south) -- node[pos=0.45, right=2pt, font=\tiny, fill=white, inner sep=1pt] {Encrypts} (f3.north);
\end{tikzpicture}
\caption{Envelope encryption architecture across three security boundaries. In Zone 1, the Master Key (KEK) is locked inside an audited hardware security module (KMS/HSM) and never touches host memory. In Zone 2, ephemeral Data Encryption Keys (DEKs) are generated on the fly to encrypt individual application objects, then wiped immediately. In Zone 3, the untrusted storage layer persists the ciphertext along with the wrapped key $\Enc_{\text{KEK}}(\text{DEK})$. A breach of persistent storage yields only encrypted keys and ciphertexts; rotation of the master key requires re-wrapping small 32-byte DEKs without rewriting gigabytes of ciphertext.}
\label{ch39:fig-envelope}
\end{figure}

Why bother? Three concrete wins. \emph{Blast radius:} the bare key that ever
touches your application memory is a single-file DEK; compromise it and you lose one
file, not the archive. \emph{Cheap rotation:} to rotate the master key you re-wrap
the small DEKs under the new KEK without re-encrypting terabytes of data ---
you touch kilobytes of wrapped keys, not the data itself. \emph{Separation of duty:}
the KMS enforces who may unwrap, with an audit log, while your app never sees the
master key at all.

\begin{example}[Envelope encryption with tiny numbers]
\label{ch39:envelope-ex}
We will use a toy XOR ``cipher'' so the arithmetic is visible; a real system uses
AES-GCM throughout. Let the master key be $K = \texttt{0x5C}$ (one byte). To store
the message byte $m = \texttt{0x41}$ (\texttt{'A'}):
\begin{enumerate}
  \item Draw a fresh DEK from the CSPRNG: $D = \texttt{0x1B}$.
  \item Encrypt data with the DEK: $c = m \xor D = \texttt{0x41} \xor \texttt{0x1B}
        = \texttt{0x5A}$.
  \item Wrap the DEK under the master key: $w = D \xor K = \texttt{0x1B} \xor
        \texttt{0x5C} = \texttt{0x47}$.
  \item Wipe $D$. Store the pair $(c, w) = (\texttt{0x5A}, \texttt{0x47})$.
\end{enumerate}
To read back: unwrap $D = w \xor K = \texttt{0x47} \xor \texttt{0x5C} =
\texttt{0x1B}$, then decrypt $m = c \xor D = \texttt{0x5A} \xor \texttt{0x1B} =
\texttt{0x41} = \text{`A'}$. Now rotate the master key to $K' = \texttt{0x33}$. We
do \emph{not} touch $c$. We only re-wrap: unwrap with the old $K$ to recover
$D = \texttt{0x1B}$, then re-wrap $w' = D \xor K' = \texttt{0x1B} \xor \texttt{0x33}
= \texttt{0x28}$. Store $(\texttt{0x5A}, \texttt{0x28})$. The 1-byte ciphertext was
never rewritten; only the 1-byte wrapped key changed. Scale $c$ to a terabyte and
this approach avoids re-encrypting the bulk ciphertext during master key rotation.
\end{example}

\begin{focusbox}
\textbf{Toy warning, sharpened.} The XOR above is for arithmetic visibility only ---
XOR-wrapping is a one-time pad and reusing $K$ across many DEKs would leak their
XORs exactly as in Example~\ref{ch39:nonce-reuse-ex}. Real envelope encryption
wraps with an AEAD (AES-GCM or AES-KW, the NIST key-wrap mode), which authenticates
the wrap and is safe under reuse of the KEK across many DEKs. Never wrap keys by
XOR. The structure of the hierarchy is the lesson; the primitive must be a real
AEAD.
\end{focusbox}

\subsection{Rotation, staged so nothing breaks}

Rotation is where good intentions cause outages. Swap the key in one atomic step
and every piece of data encrypted under the old key suddenly will not decrypt.
The fix is to let old and new keys coexist, identified by a key-id carried in the
ciphertext header --- exactly Tink's multi-key keyset (Figure~\ref{ch39:fig-rotation}).

\begin{figure}[H]
\centering
\begin{tikzpicture}[font=\small,
  st/.style={draw=sourceblue!80, fill=sourceblue!12, rounded corners=2pt, minimum width=2.7cm, minimum height=2.4cm, align=left, inner sep=4pt},
  ar/.style={-{Stealth[length=2.5mm]}, thick, structurepurple}
]
  \node[st, draw=targetgreen!80, fill=targetgreen!12] (s1) at (0, 0) {
    \textbf{1. Add Stage} \\
    \tiny $\bullet$ Generate key $v_2$ in Keyset.\\
    \tiny $\bullet$ Status: \textbf{Inactive for Enc}.\\
    \tiny $\bullet$ Status: \textbf{Active for Dec}.\\
    \tiny $\bullet$ Primary remains $v_1$.\\
    \tiny $\bullet$ All nodes sync keyset.
  };

  \node[st, draw=quotientorange!80, fill=quotientorange!12] (s2) at (3.5, 0) {
    \textbf{2. Promote Stage} \\
    \tiny $\bullet$ Mark $v_2$ as \textbf{Primary}.\\
    \tiny $\bullet$ All new writes use $v_2$.\\
    \tiny $\bullet$ Old data reads $v_1$.\\
    \tiny $\bullet$ Headers carry Key-ID:\\
    \tiny \quad $[\text{kid}: 2 \mid C_{v2}]$.\\
    \tiny $\bullet$ Zero read downtime.
  };

  \node[st, draw=structurepurple!80, fill=structurepurple!12] (s3) at (7.0, 0) {
    \textbf{3. Re-encrypt Stage} \\
    \tiny $\bullet$ Background job scans\\
    \tiny \quad persistent storage.\\
    \tiny $\bullet$ Re-wraps DEKs or\\
    \tiny \quad decrypts/re-encrypts\\
    \tiny \quad legacy records from $v_1 \to v_2$.\\
    \tiny $\bullet$ Validates migration.
  };

  \node[st, draw=trapred!80, fill=trapred!12] (s4) at (10.5, 0) {
    \textbf{4. Retire Stage} \\
    \tiny $\bullet$ Verify count of $v_1$\\
    \tiny \quad ciphertexts is 0.\\
    \tiny $\bullet$ Revoke key $v_1$.\\
    \tiny $\bullet$ Securely zeroize /\\
    \tiny \quad destroy $v_1$ in KMS.\\
    \tiny $\bullet$ Key material gone.
  };

  \draw[ar] (s1) -- (s2);
  \draw[ar] (s2) -- (s3);
  \draw[ar] (s3) -- (s4);

  \node[draw=gray!80, fill=gray!8, rounded corners=2pt, font=\scriptsize, align=center, inner sep=3pt] at (5.25, -1.8) {
    \textbf{Core Operational Invariant:} A new key can only encrypt once all nodes can decrypt it ($1 \to 2$);\\
    an old key can only be retired once zero stored ciphertexts reference it ($3 \to 4$).
  };
\end{tikzpicture}
\caption{Staged key rotation timeline. Rotating a cryptographic key in production requires maintaining backward decryption compatibility while shifting write traffic. In Stage 1, key $v_2$ is distributed to all decrypting nodes before use. In Stage 2, $v_2$ is promoted to primary for new encryptions while $v_1$ continues decrypting historical records via ciphertext key-IDs. In Stage 3, a background process re-wraps legacy records. In Stage 4, once zero active records rely on $v_1$, it is safely retired and erased.}
\label{ch39:fig-rotation}
\end{figure}

The invariant that keeps the system online: \textbf{you may start encrypting under
a new key only after every consumer can decrypt it, and you may remove an old key
only after no ciphertext still references it.} Each ciphertext names its key-id, so
the decryptor always knows which key to reach for. Skip the staging --- promote and
retire in one move --- and you have manufactured an outage or, worse, a silent data
loss when the old key is gone and old records will not open.

\subsection{Memory Hygiene: Dead-Store Elimination and Zeroization}
\label{ch39:memory-hygiene}

When a cryptographic function finishes, its raw keys, intermediate expanded states,
and plaintexts linger in heap buffers or stack frames. If those memory pages are
later reallocated without clearing, leaked via core dumps, inspected by debugging
tools, or swapped to an unencrypted disk, the key material is exposed long after
the operation finished.

The naive defense in C or C++ is:
\begin{lstlisting}[language=C]
void decrypt_and_process(const uint8_t *ciphertext, size_t len) {
    uint8_t key[32];
    derive_key(key);
    aes_gcm_decrypt(ciphertext, len, key);
    memset(key, 0, sizeof(key)); // DANGER: Compiler optimizer deletes this!
}
\end{lstlisting}

\begin{trapbox}
\textbf{The Dead-Store Elimination (DSE) trap.} Under standard compiler optimizations
(\texttt{-O2}, \texttt{-O3}), GCC and Clang analyze data dependencies. The compiler
observes that the array \texttt{key} is an automatic stack variable that is never
read again after \texttt{memset} before the function returns. Under the C ``as-if''
rule, writes to variables that are never subsequently read have no observable
effect on program behavior. The optimizer silently removes the \texttt{memset} call
entirely. The compiled binary leaves the 256-bit key sitting naked on the stack frame.
\end{trapbox}

To enforce guaranteed zeroization that survives aggressive optimization:
\begin{itemize}
  \item \textbf{C11 Annex K:} \texttt{memset\_s(key, sizeof(key), 0, sizeof(key))},
  which the ISO C standard mandates must never be optimized away.
  \item \textbf{POSIX / BSD / Linux glibc 2.25+:} \texttt{explicit\_bzero(key, sizeof(key))}.
  \item \textbf{OpenSSL / BoringSSL:} \texttt{OPENSSL\_cleanse(key, sizeof(key))},
  which uses volatile pointer arithmetic or compiler barriers to prevent DSE.
  \item \textbf{Windows API:} \texttt{SecureZeroMemory(key, sizeof(key))}.
\end{itemize}

\begin{focusbox}
\textbf{Defense in depth: memory locking with \texttt{mlock}.} Operating systems
relieve physical RAM pressure by paging memory blocks to swap partitions on disk.
If a page holding cryptographic keys is swapped out, plaintext key material is
written to non-volatile storage, surviving host reboot. High-assurance cryptographic
services call \texttt{mlock()} / \texttt{VirtualLock()} on memory regions holding
keys to instruct the kernel virtual memory manager never to write those pages to disk.
Combine with disabling core dumps (\texttt{setrlimit(RLIMIT\_CORE, 0)}) to ensure
crash dumps do not persist transient keys to disk.
\end{focusbox}

\begin{applicationbox}
\textbf{How the cloud KMSs do it.} AWS KMS, Google Cloud KMS, and HashiCorp Vault
all implement exactly this. You get a key with versions; ``rotate'' adds a new
version and makes it primary; ciphertext blobs embed the version id so old data
keeps decrypting indefinitely; the master key never leaves the service boundary
(often a FIPS-validated HSM). Envelope encryption is the default API:
\texttt{GenerateDataKey} hands you a fresh DEK \emph{and} its wrapped form in one
call, and you are expected to throw away the plaintext DEK the moment you finish
the operation. The entire design is the key hierarchy of
Example~\ref{ch39:envelope-ex}, productized.
\end{applicationbox}

\section{Secure Defaults and Crypto-Agility}
\label{ch39:defaults}

A library's defaults are a security decision made on behalf of every user who never
changes them --- which is most of them. The discipline is: \textbf{the default
configuration should be the secure configuration}, and the insecure options should
be hard to reach or absent.

This is why modern designs delete the algorithm menu. Old TLS let the client and
server \emph{negotiate} a cipher suite from a long list including export-grade
40-bit keys and broken modes; the negotiation itself became the attack surface
(FREAK, Logjam, and the downgrade attacks of Chapters~16 and~25). TLS 1.3 cut the list to a handful of AEAD suites and removed negotiable
weak options entirely. There is no insecure choice to be tricked into, because the
insecure choices are gone.

\begin{trapbox}
\textbf{Agility is a liability if it includes the past.} ``Crypto-agility'' --- the
ability to swap algorithms --- is genuinely necessary; you will need to migrate to
post-quantum primitives this decade. But agility that keeps the old algorithm
\emph{selectable} is a downgrade waiting to happen. The right shape is a versioned
\emph{suite}: version 1 is ``X25519 + AES-GCM,'' version 2 is ``ML-KEM + AES-GCM,''
and a peer either speaks version 2 or it does not --- there is no mix-and-match menu
to negotiate downward. The keyset/suite carries the choice; the wire protocol does
not haggle over primitives.
\end{trapbox}

\begin{intuitionbox}
A secure default is a gift to the 90\% of users who copy your example code verbatim
and ship it. They will not read the security section. They will not tune the
parameters. Whatever your constructor does with no arguments is what runs in
production. Treat the zero-argument path as your most important security feature,
because it is the one almost everyone actually uses.
\end{intuitionbox}

\section{Testing Crypto: Vectors, Properties, and Failure Paths}
\label{ch39:testing}

We said correctness is not testable by the implementer. That is true for
\emph{security}, but ordinary correctness --- ``does this code compute AES the way
the standard says'' --- is very testable, and you must, because a one-bit
implementation error can void all the security. The tools are test vectors,
round-trip and property tests, and --- the one people forget --- tests of the
\emph{rejection} path.

\subsection{Known-answer test vectors}

Standards bodies publish \emph{test vectors}: fixed inputs with their exact
expected outputs, computed by the reference implementation. NIST's CAVP/ACVP
vectors, Project Wycheproof from Google, RFC appendices --- they exist so your
implementation, or your wiring of a library, can be checked against ground truth.
A passing vector test proves you compute the same function the rest of the world
computes.

\begin{implementbox}
\textbf{Learning only.} This shows the \emph{shape} of vector testing, not a way to
trust homemade crypto. Check a real implementation against published vectors; never
ship the toy.
\begin{lstlisting}[language=Python, caption={Running a known-answer vector. Real
suites have thousands of these, including deliberately malformed inputs.}]
# One ChaCha20-Poly1305 test vector (from RFC 8439). The hex strings below
# are TRUNCATED for the page -- the real vector has the full 32-byte key and
# the full expected ct+tag. A vector is:
#   key, nonce, aad, plaintext -> expected ciphertext+tag.
vectors = [
    {
      "key":   bytes.fromhex("808182838485868788898a8b..."),  # full = 32 bytes
      "nonce": bytes.fromhex("070000004041424344454647"),     # 12 bytes
      "aad":   bytes.fromhex("50515253c0c1c2c3c4c5c6c7"),
      "pt":    b"Ladies and Gentlemen of the class of '99...",
      "ct_tag": bytes.fromhex("d31a8d34648e60db7b86afbc..."),  # full expected
    },
    # ... thousands more, including malformed inputs that MUST be rejected ...
]

def run_vectors(aead_encrypt, aead_decrypt, vectors):
    for v in vectors:
        out = aead_encrypt(v["key"], v["nonce"], v["aad"], v["pt"])
        assert out == v["ct_tag"], "encryption disagrees with the standard"
        back = aead_decrypt(v["key"], v["nonce"], v["aad"], v["ct_tag"])
        assert back == v["pt"],    "decryption disagrees with the standard"
\end{lstlisting}
\end{implementbox}

\subsection{Property tests and the rejection path}

Beyond fixed vectors, two kinds of generated tests catch whole classes of bug.

\emph{Round-trip property:} for random keys, nonces, and messages,
$\Dec_k(\Enc_k(m)) = m$. Easy, and it catches encoding bugs and length-handling
errors. But notice what it does \emph{not} catch: a decrypt routine that skips tag
verification round-trips perfectly. Round-trip success is necessary, never
sufficient.

\emph{Rejection property:} this is the one that matters and the one people skip.
Take a valid ciphertext, flip one bit anywhere --- in the body, in the nonce, in
the tag, in the associated data --- and assert that decryption \emph{fails}. A
shocking number of crypto bugs are decryptors that do not actually check the thing
they claim to check. The only way to catch ``the signature is never really
verified'' is a test that feeds a \emph{bad} signature and demands rejection.

\begin{attackbox}
\textbf{The verify that never verified.} A recurring real-world bug: a signature- or
MAC-checking routine that returns ``valid'' on every input because of an inverted
boolean, an exception swallowed by a bare \texttt{except: pass}, or a function that
computes the expected tag and then forgets to compare it. Apple's 2014
``goto fail'' was a sibling of this --- a duplicated \texttt{goto} skipped the final
certificate check, so TLS accepted forged certificates. Every round-trip test
passed; legitimate traffic worked perfectly. Only a \emph{negative} test --- feed a
forged certificate, assert it is rejected --- would have caught it. Always test that
bad inputs fail, not only that good inputs succeed.
\end{attackbox}

\begin{trapbox}
\textbf{Constant-time comparison, or you test your way to a timing oracle.} When you
do compare a tag or MAC, you must compare in \emph{constant time} --- a routine that
inspects every byte regardless of where the first mismatch is. The natural
\texttt{a == b} on byte strings short-circuits at the first differing byte, so its
running time leaks how many leading bytes matched. An attacker who can time
verification recovers a valid tag byte by byte, turning a forgery from $2^{128}$
guesses into a few thousand timed queries (Chapter~37 shows the full
attack). Use the library's \texttt{constant\_time\_compare} /
\texttt{sodium\_memcmp} / \texttt{hmac.compare\_digest}; never \texttt{==}.
\end{trapbox}

\begin{connectionbox}
A timing-oracle attack collapses the search cost across components. A blind forgery is a $2^{128}$ brute force --- hopeless. A
byte-at-a-time timing oracle turns it into $16$ independent searches of at most
$256$ values each --- about $128$ tries per byte on average, so roughly
$16 \times 128 \approx 2^{11}$ timed queries. The math of the
tag did not change; the comparison routine handed the attacker a way to decompose
one impossible search into sixteen trivial ones. Constant-time code denies that
decomposition.
\end{connectionbox}

\section{Constant-Time Discipline and Branch Elimination}
\label{ch39:constant-time}

Every programmer is taught to optimize code with early exits and conditional branches:
if a condition is false, skip the work; if a match is found, return immediately.
In cryptography, that instinct is lethal. The physical CPU does not compute in an
abstract mathematical ether. It executes on silicon pipelines with branch prediction
units (BPUs), instruction caches, out-of-order execution windows, and arithmetic
logic units (ALUs).

When code branches on a condition that depends on secret data --- a key bit, a
padding byte, an authentication tag --- the CPU's branch predictor records the
outcome in internal history tables. An adversary sharing the CPU core (e.g., in a
cloud environment) or timing network packet responses can observe whether the branch
was taken.

\subsection{The Constant-Time Multiplexer}
\label{ch39:ct-mux}

The fundamental building block of branchless programming is the \emph{constant-time
conditional select} (or multiplexer). Suppose we have two $w$-bit values $a$ and $b$,
and a secret decision flag $c \in \{0, 1\}$. We wish to compute:
\[
r = \begin{cases} a & \text{if } c = 1, \\ b & \text{if } c = 0, \end{cases}
\]
without executing an \texttt{if/else} statement or conditional jump instruction.

\begin{figure}[H]
\centering
\begin{tikzpicture}
  % Panel A: The Timing Leak of Branching vs Bitwise Multiplexer
  \begin{scope}[shift={(0,0)}]
    \node[font=\small\bfseries] at (2.8, 3.8) {Panel A: Branching vs. Bitwise Multiplexer};

    % Branching code box
    \node[draw=trapred!80, fill=trapred!10, rounded corners=2pt, font=\scriptsize, align=left, inner sep=4pt] (br) at (2.8, 2.4) {
      \textbf{Vulnerable Branching:} \\
      \texttt{if (c == 1) r = a; else r = b;}\\
      $\bullet$ Emits conditional branch (\texttt{JNE} / \texttt{CBZ}).\\
      $\bullet$ BPU history table updated with secret $c$.\\
      $\bullet$ Speculative execution side channels (Spectre).\\
      $\bullet$ Variable instruction cache \& latency.
    };

    % Constant time code box
    \node[draw=targetgreen!80, fill=targetgreen!10, rounded corners=2pt, font=\scriptsize, align=left, inner sep=4pt] (ct) at (2.8, 0.4) {
      \textbf{Constant-Time Multiplexer:} \\
      \texttt{mask = -c;  r = b \^{} ((a \^{} b) \& mask);}\\
      $\bullet$ Zero branches, zero jump instructions.\\
      $\bullet$ Constant cycle count on all modern CPUs.\\
      $\bullet$ Arithmetic sign expansion: $-0 = \texttt{0x00}$, $-1 = \texttt{0xFF}$.\\
      $\bullet$ Fully immune to branch-prediction attacks.
    };
  \end{scope}

  % Panel B: Bit-Level Circuit Dataflow
  \begin{scope}[shift={(6.4,0)}]
    \node[font=\small\bfseries] at (3.1, 3.8) {Panel B: Bit-Level Logic Dataflow};

    % Inputs
    \node[draw=sourceblue!80, fill=sourceblue!15, rounded corners=2pt, font=\scriptsize\bfseries, minimum width=1.4cm] (inA) at (0.8, 2.9) {Input $a$};
    \node[draw=targetgreen!80, fill=targetgreen!15, rounded corners=2pt, font=\scriptsize\bfseries, minimum width=1.4cm] (inB) at (3.2, 2.9) {Input $b$};
    \node[draw=quotientorange!80, fill=quotientorange!15, rounded corners=2pt, font=\scriptsize\bfseries, minimum width=1.8cm] (inC) at (5.4, 2.9) {Condition $c \in \{0,1\}$};

    % Two's complement mask expansion
    \node[draw=structurepurple!80, fill=structurepurple!15, rounded corners=2pt, font=\scriptsize, align=center, inner sep=2pt] (mask) at (5.4, 1.8) {
      $\text{mask} = -c$\\ \tiny $-0 = \texttt{00\dots0}$\\ \tiny $-1 = \texttt{FF\dots F}$
    };
    \draw[-{Stealth[length=2mm]}, thick] (inC) -- (mask);

    % XOR 1: (a ^ b)
    \node[draw=structurepurple!80, circle, inner sep=1pt, fill=white, font=\scriptsize\bfseries] (xor1) at (2.0, 1.8) {$\oplus$};
    \draw[-{Stealth[length=2mm]}, thick] (inA) -- (xor1);
    \draw[-{Stealth[length=2mm]}, thick] (inB) -- (xor1);
    \node[font=\tiny, left=1pt] at (xor1) {$a \oplus b$};

    % AND gate: ((a ^ b) & mask)
    \node[draw=structurepurple!80, fill=structurepurple!15, rounded corners=2pt, font=\scriptsize\bfseries, inner sep=3pt] (andg) at (3.7, 0.7) {$\&$};
    \draw[-{Stealth[length=2mm]}, thick] (xor1) |- (andg.west);
    \draw[-{Stealth[length=2mm]}, thick] (mask) |- (andg.east);

    % Final XOR: b ^ ((a ^ b) & mask)
    \node[draw=targetgreen!80, circle, inner sep=1pt, fill=white, font=\scriptsize\bfseries] (xor2) at (3.7, -0.6) {$\oplus$};
    \draw[-{Stealth[length=2mm]}, thick] (andg) -- (xor2);

    % Route b down to final XOR
    \draw[-{Stealth[length=2mm]}, thick, dashed] (inB.south) -- ++(0,-0.4) -| (1.0, 0.0) |- (xor2.west);

    % Output
    \node[draw=targetgreen!80, fill=targetgreen!20, rounded corners=2pt, font=\scriptsize\bfseries, minimum width=2.4cm] (outR) at (3.7, -1.5) {Result $r$};
    \draw[-{Stealth[length=2mm]}, thick] (xor2) -- (outR);

    % Truth table note below Result r
    \node[draw=gray!80, fill=gray!8, rounded corners=2pt, font=\tiny, align=left, inner sep=3pt] at (3.1, -2.5) {
      $c=0 \implies \text{mask}=0 \implies r = b \oplus 0 = b$\\
      $c=1 \implies \text{mask}=\texttt{\~{}0} \implies r = b \oplus (a \oplus b) = a$
    };
  \end{scope}
\end{tikzpicture}
\caption{Constant-time multiplexer bit-level logic diagram ($r = b \oplus ((a \oplus b) \ \& \ -c)$). In Panel A, conditional branching code (\texttt{if/else}) introduces variable cycle latency and alters branch predictor state, creating exploitable timing side channels. In Panel B, the selection is performed using straight-line bitwise operations. Two's complement negation converts a single bit condition $c \in \{0, 1\}$ into an all-zero or all-one bitmask ($\text{mask} = -c$). The bitwise formula computes $r = b \oplus ((a \oplus b) \ \& \ \text{mask})$, which evaluates to $b$ when $c = 0$ and $a$ when $c = 1$ in identical clock cycles on every modern CPU.}
\label{ch39:fig-ct-mux}
\end{figure}

The mechanism relies on two's complement sign expansion and bitwise Boolean logic (Figure~\ref{ch39:fig-ct-mux}):
\begin{enumerate}
  \item \textbf{Mask generation via negation:} Compute $\text{mask} = -c$. In C, write this on an \emph{unsigned} type, \texttt{mask = -(uint32\_t)c}, so the wraparound is well defined. In two's complement arithmetic:
  \[
  -0 = 0 = \texttt{0x00\dots00}, \qquad -1 = \texttt{0xFF\dots FF} \quad (\text{all bits set to } 1).
  \]
  \item \textbf{Bitwise multiplexing:} Compute:
  \[
  r = b \oplus \big((a \oplus b) \ \& \ \text{mask}\big).
  \]
\end{enumerate}

Trace both cases:
\begin{itemize}
  \item \textbf{When $c = 0$:} $\text{mask} = \texttt{0x00\dots00}$. Then $(a \oplus b) \ \& \ 0 = 0$, giving:
  \[
  r = b \oplus 0 = b.
  \]
  \item \textbf{When $c = 1$:} $\text{mask} = \texttt{0xFF\dots FF}$. Then $(a \oplus b) \ \& \ \texttt{\~{}0} = (a \oplus b)$, giving:
  \[
  r = b \oplus (a \oplus b) = a.
  \]
\end{itemize}

Every CPU instruction executed --- negation, XOR, AND, XOR --- takes fixed cycle latency regardless of whether $c$ is 0 or 1. No jump is taken; no branch predictor entry is dirtied; no speculative window is opened. One caveat: an optimizing compiler (\texttt{-O2}/\texttt{-O3}) may recognize the pattern and turn it back into a branch or a \texttt{cmov}, so constant-time libraries hide $c$ behind a value barrier (an empty inline-\texttt{asm} or a \texttt{volatile} read) and check the emitted assembly.

\section{Threat Modeling: Decide What You Are Defending Against}
\label{ch39:threat}

You cannot build a secure system without first writing down what ``secure'' means
for it --- against whom, protecting what, assuming what. Skip this and you will
either under-protect (no defense against an attacker who is clearly in scope) or
over-engineer (HSMs and key ceremonies to guard a public RSS feed). A threat model
is four questions answered in plain language before any code.

\begin{enumerate}
  \item \textbf{What are you protecting?} The assets: user messages, password
  database, signing key, session tokens. Be specific; ``the data'' is not an asset.
  \item \textbf{Who is the adversary, and what can they do?} The capabilities: a
  passive eavesdropper on the wire? An active attacker who can inject and reorder
  (Dolev--Yao)? A malicious insider with database read access? An attacker who can
  time your server? Capabilities, not personas.
  \item \textbf{What must hold for you to be secure?} The security goals tied to
  assets: confidentiality of messages, integrity of transfers, authenticity of the
  server, forward secrecy if a key later leaks.
  \item \textbf{What is explicitly \emph{out} of scope?} The assumptions you are
  making. ``We do not defend against a compromised client device.'' ``We assume the
  CSPRNG is sound.'' Writing these assumptions down explicitly separates
  documented trust boundaries from hidden blind spots.
\end{enumerate}

\begin{example}[A threat model for a small message store]
\label{ch39:threatmodel-ex}
Suppose you are building a service where users store notes and read them back later.
A one-page model:

\textbf{Assets:} the note contents (confidentiality), and the fact that a returned
note is the one the user actually stored (integrity).

\textbf{Adversaries and capabilities:} (a) a network attacker who can read and
modify traffic between client and server --- handled by TLS, in scope; (b) an
attacker who steals a database backup --- in scope, so notes must be encrypted at
rest under a key \emph{not} in that backup; (c) an attacker who fully compromises
the running server, including memory and the live KEK --- declared \emph{out} of
scope for this tier of product.

\textbf{Goals:} TLS 1.3 for the channel; per-note envelope encryption with DEKs
wrapped by a KMS-held KEK for at-rest confidentiality and integrity; the KEK lives
in the KMS so a stolen backup is useless without KMS access.

\textbf{Out of scope, stated:} a live-server compromise that can call the KMS to
unwrap; client-device malware; traffic-analysis (sizes and timing). These are
acknowledged, not defended, at this tier.

Notice the model immediately \emph{decides the architecture}: adversary (b) is the
reason DEKs are wrapped by a KMS key rather than a key stored next to the data (Figure~\ref{ch39:fig-threat-arch}).
Without writing down (b) you might have ``encrypted at rest'' with the key in the
same backup --- security theater that the model exposes in one line.
\end{example}

\begin{figure}[H]
\centering
\begin{tikzpicture}[font=\footnotesize,
  nodebox/.style={draw=sourceblue!80, fill=sourceblue!10, rounded corners=3pt, minimum width=2.7cm, minimum height=1.1cm, align=center},
  kmsbox/.style={draw=structurepurple!80, fill=structurepurple!12, rounded corners=3pt, minimum width=2.8cm, minimum height=1.1cm, align=center},
  dbbox/.style={draw=targetgreen!80!black, fill=targetgreen!12, rounded corners=3pt, minimum width=3.1cm, minimum height=1.1cm, align=center},
  adv/.style={draw=trapred!85, fill=trapred!10, rounded corners=2pt, inner sep=3.5pt, align=center, font=\scriptsize},
  arr/.style={-{Stealth[length=2.2mm]}, thick}
]
  \node[nodebox] (cli) at (0, 0) {\textbf{Client}\\User notes};
  \node[nodebox] (srv) at (4.4, 0) {\textbf{App Server}\\Generates DEK\\Encrypts note};
  \node[dbbox] (db) at (9.4, 0) {\textbf{Database \& Backups}\\$\Enc_{\text{DEK}}(\text{note})$\\$+\,\text{Wrap}_{\text{KEK}}(\text{DEK})$};
  \node[kmsbox] (kms) at (4.4, 2.1) {\textbf{KMS / HSM}\\Holds KEK (never exported)};

  \draw[{Stealth[length=2mm]}-{Stealth[length=2mm]}, very thick, sourceblue!80!black] (cli.east) -- node[above, font=\scriptsize\bfseries] {TLS 1.3} (srv.west);
  \draw[arr, sourceblue!80!black] (srv.east) -- node[above, font=\scriptsize] {Ciphertext} (db.west);
  \draw[{Stealth[length=2mm]}-{Stealth[length=2mm]}, thick, structurepurple!80!black] (srv.north) -- node[right, font=\scriptsize] {Wrap / unwrap DEK} (kms.south);

  % Adversaries
  \node[adv] (advA) at (2.2, -1.5) {\textbf{Adversary (a): Wire Attacker}\\Blocked by TLS 1.3 AEAD};
  \node[adv] (advB) at (9.4, -1.5) {\textbf{Adversary (b): Backup Thief}\\Gets wrapped DEKs, no KEK};
  \draw[-{Stealth[length=2mm]}, dashed, trapred!80, thick] (advA.north) -- (2.2, -0.05);
  \draw[-{Stealth[length=2mm]}, dashed, trapred!80, thick] (advB.north) -- (db.south);
\end{tikzpicture}
\caption{Architecture driven by the threat model of Example~\ref{ch39:threatmodel-ex}. Adversary~(a) on the network is stopped by TLS~1.3; Adversary~(b) who steals a database backup gets only ciphertexts and wrapped DEKs because the KEK remains isolated in the KMS.}
\label{ch39:fig-threat-arch}
\end{figure}

\begin{intuitionbox}
A threat model is a list with two columns: attacks you will stop, and attacks
you have decided not to. Write both down. Without the first column you leave
holes; without the second you spend months hardening against an adversary who
was never coming.
\end{intuitionbox}

\section{When (Almost Never) to Implement a Primitive}
\label{ch39:when-roll}

Having said ``almost never'' many times, intellectual integrity demands we name the
cases where touching a primitive is defensible --- and surround each with the
conditions that make it not reckless.

\begin{itemize}
  \item \textbf{You are a cryptographer doing research}, designing or analyzing a
  scheme for public scrutiny. The deliverable is a paper and a reference
  implementation meant to be attacked, not a product feature.
  \item \textbf{No vetted library exists for your platform.} A genuinely novel
  environment --- a tiny microcontroller, an exotic enclave, a new language with no
  binding to libsodium/BoringSSL --- may force your hand. The right move is usually
  to port a vetted implementation, not to write fresh code, and to get it reviewed
  and tested against full vector suites and constant-time checks before trusting it.
  \item \textbf{You are deliberately learning} --- working the examples in a book
  like this one. This is valuable and you should do it. The discipline is to
  \emph{label} the result as a learning artifact and never let it cross into a real
  system. That is exactly why every code box here carries a learning-only warning.
\end{itemize}

If you do find yourself at Layer 0 for a real deployment, the bar is high: a
public, peer-reviewed algorithm (never a secret design); constant-time code with
no secret-dependent branches or memory indices (Chapter~37); the full
published vector suite passing, including malformed-input rejection; a defended
RNG; secure memory handling that wipes keys and resists being swapped to disk; and,
non-negotiably, an independent expert audit. If you cannot fund the audit, you
cannot afford to roll the primitive. A bug in homemade crypto is silent: it
passes every functional test and fails only when an attacker finds it.

\begin{applicationbox}
The libraries you should be reaching for are themselves the product of this
discipline done right. BoringSSL is Google's audited fork of OpenSSL, born after
Heartbleed showed what a single unchecked length in a widely-trusted C codebase can
do. libsodium is a hardened, portable packaging of NaCl. \texttt{ring} and RustCrypto bring
memory-safe implementations with constant-time intent baked in. Each represents
many person-years of expert attention, fuzzing, formal verification, and audit
--- precisely the work you are \emph{not} going to replicate over a weekend, and the
reason standing on their shoulders is not laziness but engineering.
\end{applicationbox}

\section{A Practical Checklist}
\label{ch39:checklist}

Print this. Tape it next to the review queue. It compresses the chapter into
questions you can answer yes/no during a code review, grouped by the failure
classes we have walked through.

\begin{focusbox}
\textbf{Shipping-cryptography checklist.}

\emph{Library and layer}
\begin{itemize}
  \item Am I using a vetted, audited, maintained library --- not my own primitive?
  \item Am I standing on the highest layer that solves my problem (a box/AEAD API,
  not raw modes)?
  \item Is every algorithm choice a current standard (AES-GCM or ChaCha20-Poly1305
  for AEAD, HKDF for derivation, X25519/ML-KEM for exchange, Ed25519 for
  signatures, Argon2/scrypt for passwords)? No DES, RC4, MD5, SHA-1, ECB, or
  textbook RSA.
\end{itemize}

\emph{Keys}
\begin{itemize}
  \item Are keys generated from a CSPRNG with full entropy --- never a password
  directly, never a constant, never \texttt{rand()}?
  \item Are passwords run through a real KDF (Argon2id, scrypt, or PBKDF2) with a
  per-user salt --- never hashed bare?
  \item Are long-lived keys in a KMS/HSM, out of source control, out of logs, out of
  environment dumps? Is a secret-scanner in the commit hook?
  \item Is each key used for exactly one purpose? No key reuse across schemes.
  \item Is there a staged rotation plan (add, promote, re-encrypt, retire) with
  key-ids in the ciphertext?
\end{itemize}

\emph{Nonces, modes, and composition}
\begin{itemize}
  \item Are nonces/IVs unique per (key, message) --- counter-managed for 96-bit
  GCM, random only with extended nonces? Survives restarts, forks, and VM
  snapshots?
  \item Is encryption authenticated (AEAD)? Never confidentiality without
  integrity.
  \item If composing by hand (you should not be), is it encrypt-then-MAC, never
  MAC-then-encrypt or encrypt-and-MAC?
\end{itemize}

\emph{Comparison, randomness, and side channels}
\begin{itemize}
  \item Are all secret comparisons constant-time
  (\texttt{compare\_digest}/\texttt{sodium\_memcmp})? Never \texttt{==} on a tag.
  \item Is there no secret-dependent branch or memory index in code you wrote?
  \item Is the RNG the OS CSPRNG, seeded, never reused after fork?
\end{itemize}

\emph{Defaults, testing, and modeling}
\begin{itemize}
  \item Is the zero-argument / default path the secure path?
  \item Do tests run published vectors \emph{and} negative tests (tampered input
  MUST be rejected; a forged signature MUST fail)?
  \item Is there a written threat model naming assets, adversary capabilities,
  goals, and explicit out-of-scope assumptions?
  \item Is there a plan for crypto-agility (versioned suites, no downgrade menu)
  toward post-quantum?
\end{itemize}
\end{focusbox}

\begin{applicationbox}
\textbf{What the systems you trust actually do.} Signal stands on libsignal (a
vetted protocol implementation) rather than reinventing the ratchet per platform.
WireGuard ships a single fixed cipher suite --- Curve25519, ChaCha20-Poly1305,
BLAKE2s --- with no negotiation, so there is no downgrade to attack; ``no algorithm
agility on the wire'' is a stated design goal. Apple and Google's mobile keystores
keep device keys in dedicated secure hardware (Secure Enclave, Titan/StrongBox), so
even a fully compromised OS cannot export them --- the key-hierarchy and
out-of-scope-compromise ideas from our threat model, in silicon. age and Tink give
application developers misuse-resistant file and data encryption with managed keys.
Across all of them the through-line is identical to this chapter: \emph{choose
strong primitives once, expose a narrow safe interface, and refuse to let the caller
make the dangerous decision.}
\end{applicationbox}

\section{What to Carry Forward}
\label{ch39:carry-forward}

The mathematics in this book is the reason cryptography \emph{can} work. This
chapter is the reason it sometimes does not. Keep the central asymmetry in view:
the cipher almost never fails, whereas the wiring around it fails silently whenever
an adversary probes an overlooked assumption. Reliable cryptographic engineering
means putting fewer decisions in application code and delegating them to libraries
that have already survived sustained review.

Stand as high on the API ladder as your problem allows. Reach for libsodium or
Tink, or higher still for a whole protocol you do not assemble. Let the library
own the nonce, the mode, the comparison, the composition. Manage keys with a
hierarchy and a KMS, rotate in stages, and keep secrets out of the repo. Make the
default the secure default. Test the rejection path, not only the round trip. Write
the threat model before the code. And treat ``implement a primitive yourself'' as a
decision that needs an audit budget, not a free afternoon. Do that, and the
mathematics of the preceding chapters gets to do its job ---
which is, in the end, the only job: to keep the secret.

To keep all the supporting mathematics in one place for quick lookup, the first appendix gathers the foundational number theory --- from the Euclidean algorithm and B{\'e}zout's identity to Euler's totient theorem and primality testing --- that underpins public-key cryptography.

\section*{Exercises}
\addcontentsline{toc}{section}{Exercises}

\begin{exercise}[Nonce reuse by hand]
A CTR-based AEAD produces keystream byte $S = \texttt{0x6D}$ for a given
(key, nonce). Two messages are sent under the same key and the \emph{same} nonce:
$m_1 = \text{`}\texttt{T}\text{'} = \texttt{0x54}$ and
$m_2 = \text{`}\texttt{F}\text{'} = \texttt{0x46}$ (a true/false bit). Compute
$c_1$, $c_2$, and $c_1 \xor c_2$. Verify $c_1 \xor c_2 = m_1 \xor m_2$ and explain,
in one sentence, what Eve learns without ever recovering $S$.
\end{exercise}

\begin{exercise}[Envelope re-wrap]
Using the toy XOR scheme of Example~\ref{ch39:envelope-ex}: master key
$K = \texttt{0x2A}$, data byte $m = \texttt{0x6B}$, fresh DEK $D = \texttt{0x39}$.
(a) Compute the stored pair $(c, w)$. (b) The master key rotates to
$K' = \texttt{0x91}$. Compute the new wrapped key $w'$ \emph{without} changing $c$.
(c) State the one quantity that gets rewritten on a master-key rotation and the one
that does not, and explain why this is the central payoff of envelope encryption.
\end{exercise}

\begin{exercise}[Count the rope]
For each API signature, list the security-relevant decision(s) the caller is forced
to make, and rank the three from most to least misuse-resistant:\\
(a) \texttt{secretbox(message, key)};\\
(b) \texttt{aes\_gcm\_encrypt(message,}\allowbreak \texttt{key, nonce, aad)};\\
(c) \texttt{aes\_cbc\_encrypt(message,}\allowbreak \texttt{key, iv)} followed by a
separate \texttt{hmac(...)} you must order and compare yourself. Justify the ranking by the
``count the arguments'' heuristic and name the specific catastrophic failure each
delegated decision enables.
\end{exercise}

\begin{exercise}[The verify that never verified]
Here is a sketch of a MAC check:
\begin{quote}\ttfamily
def verify(key, msg, tag):\\
\hspace*{2em}expected = hmac(key, msg)\\
\hspace*{2em}return True
\end{quote}
(a) Which test --- round-trip or rejection --- catches this bug, and which one
passes despite it? (b) Write, in words, the single negative test case that exposes
it. (c) Relate this to Apple's ``goto fail'': what was structurally the same?
\end{exercise}

\begin{exercise}[Timing oracle arithmetic]
A 16-byte authentication tag is verified with a comparison that short-circuits at
the first differing byte, and an attacker can measure the verification time. (a)
Estimate the number of timed queries to recover the full tag byte by byte, and
compare it to the $2^{128}$ cost of a blind forgery. (b) Name the one-line code
change that defeats the attack and explain why it removes the leak. (c) Connect
this to the search-cost-collapse idea from the connection box.
\end{exercise}

\begin{exercise}[Write a threat model]
You are building a password manager: users store credentials, encrypted, synced
across their devices through your server. Write a one-page threat model in the
four-question form of Example~\ref{ch39:threatmodel-ex}. Be explicit about whether
a \emph{server compromise} is in scope, and let that decision drive where the
master key lives and what the server can and cannot decrypt. Identify at least one
adversary you place \emph{out} of scope and say why.
\end{exercise}

\begin{exercise}[Staged rotation invariant]
A team rotates an encryption key by deleting the old key and installing the new one
in a single deploy. (a) Describe precisely what happens to data encrypted under the
old key, and what the user experiences. (b) State the two-part safety invariant
from Section~\ref{ch39:keymgmt} that the team violated. (c) Give the corrected
four-phase sequence and explain what the key-id in each ciphertext header is for.
\end{exercise}

\begin{exercise}[Defaults as a security feature]
Library A's AEAD constructor with no arguments selects AES-GCM and generates a
random 96-bit nonce per message from the OS CSPRNG. Library B's selects
XSalsa20-Poly1305 with a 192-bit random nonce. (a) Using the birthday bound from
the connection box, estimate how many messages each can safely encrypt before nonce
collision becomes a realistic worry. (b) Argue which default is more
misuse-resistant for a developer who never reads the docs, and tie your answer to
the ``zero-argument path'' principle.
\end{exercise}

\begin{exercise}[Justify the exception]
Give a concrete scenario in which implementing a standard primitive yourself
(Layer~0) is defensible, then list every condition from Section~\ref{ch39:when-roll}
that must hold before the result may touch a real system. For one of those
conditions --- constant-time code --- give an example of a seemingly innocent line
that violates it (a secret-dependent branch or a table lookup indexed by secret
data) and explain the leak.
\end{exercise}

\begin{exercise}[Audit a snippet, end to end]
You are reviewing a service that: (i) reads a 16-byte AES key from an environment
variable set in the deploy YAML committed to git; (ii) encrypts records with
AES-CBC and a fixed all-zero IV; (iii) stores the ciphertext with no MAC; (iv)
verifies an API token with \texttt{if token == stored\_token}. Identify every
distinct vulnerability against the checklist in Section~\ref{ch39:checklist}, and
for each give the one-line remediation. (There are at least five.)
\end{exercise}
